\documentclass[12pt,a4paper,onecolumn]{article}
\title{Digest of The New York Times}
\author{The New York Times}
% \definecolor{dark-green}{rgb}{0,0.25,0}
% \usepackage[utf8]{inputenc}

\usepackage{fontspec,xunicode,xltxtra}
\setmainfont[Ligatures=Required,Mapping=tex-text]{Georgia}
% \setmainfont[Ligatures=Required,Mapping=tex-text]{Palatino Linotype}
\setsansfont[Ligatures=Required,Mapping=tex-text]{Myriad Pro}
\setmonofont[Ligatures=Required,Mapping=tex-text]{Courier New}
\usepackage{xeCJK}
\setCJKmainfont{宋体}
\setCJKmonofont{黑体}

% \usepackage{indentfirst}


\usepackage[dvips,dvipsnames,svgnames]{xcolor}
\definecolor{light-gray}{gray}{0.95}

%生成PDF的链接
\usepackage[dvipdfmx]{hyperref}
\hypersetup{
    bookmarks=true,         % show bookmarks bar?
    bookmarksopen=true,
    pdfpagemode=UseNone,    % options: UseNode, UseThumbs, UseOutlines, FullScreen
    pdfstartview=FitH,
    pdfborder=1,
    pdfhighlight=/P,
    pdfauthor={wuxch},
    unicode=true,           % non-Latin characters in Acrobat’s bookmarks
    colorlinks,             % false: boxed links; true: colored links
    linkcolor=blue,         % color of internal links
    citecolor=blue,        % color of links to bibliography
    filecolor=magenta,      % color of file links
    urlcolor=cyan           % color of external links
}
% 使用目录
\usepackage{makeidx}
% 使用图形
\usepackage{graphicx}


\usepackage{titletoc}
\dottedcontents{section}[1.5em]{}{2.3em}{1pc}

% % 页眉页脚的设定
% \usepackage{fancyhdr}
% \pagestyle{fancy}
% \lhead{\includegraphics[height=0.2in]{zte_logo}}

% \usepackage{footmisc}
\usepackage{pfnote}
\usepackage{fnpos}

% % 嵌入的代码显示
% \usepackage{listings}
% \lstset{language=C++, breaklines, extendedchars=false}
% \lstset{keywordstyle=\color{blue},
%         commentstyle=\color{SeaGreen},
%         stringstyle=\ttfamily,
%         showstringspaces=false,
%         backgroundcolor=\color{light-gray}}
% \lstset{basicstyle=\small}

\usepackage{booktabs}

\usepackage{tipa}
\usepackage{phonetic}
\usepackage[gen]{eurosym}

% \usepackage{CJK}
\usepackage{makeidx}
% \usepackage{glossary}

\usepackage{fancyhdr}

\usepackage{lastpage}
\cfoot{\textsf{\thepage\ of \pageref*{LastPage}}}

\pagestyle{fancy}
\fancyhead{}
\renewcommand{\headrulewidth}{0pt}
\pagenumbering{arabic}

\usepackage{graphicx}
\usepackage{picinpar}

\setlength{\parindent}{0in}
\setlength{\parskip}{2.5ex plus4mm minus3mm}
% \setlength{\parskip}{.3\baselineskip}

\setcounter{section}{0}

% bibtex
\usepackage[super]{natbib}
% \usepackage[globalcitecopy,labelstoglobalaux]{bibunits}
% \makeatletter
% \renewcommand{\@bibunitname}{\jobname.\the\@bibunitauxcnt}
% \makeatother

% \renewcommand{\refname}{缩略语}
% \renewcommand{\refname}{}

\makeatletter                   % Reference list option change
\renewcommand\@biblabel[1]{#1:  } % from [1] to 1: 
\makeatother                    %

\newcommand{\citenumfont}[1]{\textsf{#1}}

% \usepackage{tipa}
% \usepackage{phonetic}

\newcommand{\phonetic}[1]{{\fontspec[Colour=CC00CC]{Microsoft Sans Serif} #1}}


%%% Local Variables:
%%% TeX-master: "medical_trivia"
%%% End:

\makeindex
\begin{document}
\date{}

% \maketitle
\tableofcontents
\pagebreak
\section{New Inks Cut Costs of Office Color Printing}

Xerox thinks it has hit on something big to breathe life into the ailing market for high-volume
office printers: It's called color.

After years of research, Xerox will release the first in a new series of large machines this month
that it claims will change the economics of printing large volumes of color documents at offices.
Unlike traditional laser printers that use cartridges of powdery, sometimes messy toner, the
revamped products rely on hunks of ink that remain solid at room temperature and then melt when
heated.

Xerox, based in Norwalk, Conn., bills the solid ink technology, which lets it create machines with
fewer parts, as the biggest shift in the industry since it sold the first laser printer in 1977.

``This is something they have been working on for forever and a day,'' said Brian Bissett, the
editor of the MFP Report, a publication about office printers and copiers. ``It does have the
potential to be fairly important.''

Traditionally, printer makers have charged big corporate customers a per-page fee for printing
services: up to 8 cents a page for a color document and less than 2 cents a page for
black-and-white.

Xerox and its main competitors, Canon and Ricoh, typically bill customers for a color page if there
is any pigment on it, whether it's a full-color replica of a slide from a presentation or a page
that just has a colorful company logo and mostly black text.

Because of the higher cost, only about 15 percent of the more than two trillion documents printed
every year in offices worldwide use color.

``This isn't because people don't like color but because there have been lots of barriers to
color,'' said Ursula M.~Burns, the president of Xerox. ``The last barrier is that it's expensive.''

With its new ColorQube solid ink systems, Xerox says it can pull the average cost of color documents
down to about 3 cents a page. A portion of the price break comes from better tracking of how much
color is used on a page. But the principal breakthrough with the solid ink technology stems from the
overall simplicity of the machine's innards. Solid ink machines need about one-third as many parts
as laser printers, which require various mechanisms to fuse toner onto paper that must be replaced
over time, adding materials and services costs.

And instead of putting clunky cartridges of toner into the printers, customers just pop in the
sticks of ink. Once that ink runs out, another stick goes in, and there's nothing that needs to be
thrown away or recycled. Over all, Xerox claims that it takes about 272 pounds of material to
produce one million pages with a solid ink machine versus 965 pounds of material with a competing
laser product.

Hewlett-Packard has tried to disrupt the high-volume printer market with a different type of
innovation: its Edgeline machines run on liquid ink instead of toner. Software and hardware tracking
systems within the Edgeline products also make it possible to offer customers different prices per
page based on how much color ink is actually consumed.

However, analysts say that H.P. is at a disadvantage because the technology is unfamiliar and
because H.P. lacks the sales reach of Xerox in this category.

\section{For Small Employers, Shedding Workers and Tears}

GRAFTON, Wis. -- In a small, windowless conference room, the nine members of the management team at
Ram Tool gathered to consider which employees should be laid off in the company's latest round of
cutbacks. They debated each name and weighed issues like seniority and skills. Could they do
multiple jobs? What was their attendance record?

Finally, after three days of discussions, they arrived at a list, and it fell to Shelly Polum, the
vice president for administration at this small, family-owned tool-and-die manufacturing company
here, to inform four workers they were being let go. She put on what her husband called her
``stone-cold face'' and walked out onto the shop floor.

When it was over, trying to maintain her composure, she rushed back to her office and shut the door
quickly.

Then she sank to the floor and burst into tears.

``It was a rush of emotions,'' Mrs.~Polum said later. ``I think what really hit me was, I wanted
this to be the last time. Is this enough? Is this the end? I guess I hoped this was the last time I
have to do this.''

Since the economy went into a freefall last September, Mrs.~Polum, 43, has had the task of laying
off 26 workers. It is a delicate and often painful assignment being carried out across the country
by business owners, human resources executives and other managers, and it can carry a large
emotional toll. That is particularly true at small and medium-size companies with fewer than 500
employees, like Ram Tool, which employ half the nation's private sector workers and where cutbacks
are often very personal.

Charlie Thomas III, vice president of Shuqualak Lumber in Shuqualak, Miss., had to dismiss nearly a
quarter of his work force, which used to number about 160 employees, at the end of October. He also
dismissed a handful of workers in January.

He wrote out a speech for the announcement in October in front of his men, whom he told to gather in
the lumberyard. Midway through delivering it, Mr.~Thomas had to stop and go back into his office to
compose himself.

``I couldn't get it out,'' he said. ``It just killed my soul.''

Ram Tool, which was founded 30 years ago by Mrs.~Polum's father, Roy Kannenberg, in his garage, had
nearly 100 employees as recently as several years ago but is now down to about 50.

The company's customers come from industries like oil and gas, construction and automotive, all
sectors hurt by the economic downturn. Sales, which began to slow in late summer, have plummeted by
half, leaving the company struggling for survival. After slashing hours and making other cutbacks,
the company began the process of laying off workers several at a time in October.

In the management team meetings, each member, armed with an employee roster by department, suggested
workers to be cut. A vote was taken on each worker. Some got reprieves; others were added to the
list.

A premium was put on workers with diversified skills. One employee was saved because he happened to
have a commercial driver's license, so he could make deliveries as well.

There were impossible dilemmas: what to do about two equally valued employees when one was single
and the other had a family to support; how to balance an employee's talent against his cost.

Once an initial list of names was settled upon, the small group of family members who own and run
the company met again to finalize the decision. Over the last few months, they have even approved
the layoffs of two of their relatives.

But it has been left to Mrs.~Polum to carry out the decisions. Some of her relatives have offered to
relieve her of the duty, but she insists that it is part of her job.

``I wanted to make sure I knew what was being said to that employee,'' she said. ``I wanted to be
the person that said it, and I wanted to be the last person that had anything to say.''

The process of steeling herself begins the night before, when Mrs.~Polum said she usually cannot
bring herself to eat.

``It affects me so internally because I know I'm going to hurt somebody,'' she said. ``I'm affecting
their lives.''

On their half-hour drive into work from their home in Slinger, Wis., Mrs.~Polum's husband, Dennis,
who handles purchasing, said he instructs her to block out her feelings.

``I call it a `stone-cold face,' '' Mr.~Polum said. ``You can't show any emotions. You almost have
to step outside of yourself to do it.''

On the morning of the most recent layoff in late March, Mrs.~Polum called an employee meeting in the
lunchroom. She explained the company's deteriorating financial situation and informed them some
people would have to be let go.

Instead of calling workers to her office, as she had done in the past, Mrs.~Polum has been pulling
workers aside after handing them their weekly paycheck.

Just before she went out on the floor, Mrs.~Polum wrote on the palm of her hand the names of the
workers she would have to approach.

``To make sure I'm hitting the right people,'' she explained.

It was when Mrs.~Polum got to the final worker on her palm that her emotions finally caught up to
her. The employee and his wife care for several adopted children. Even though his department had
been slow for months, he had been deliberately spared by the managers in previous layoffs.

``We were all avoiding it,'' said Mike Kannenberg, Mrs.~Polum's brother, who oversees manufacturing.

But now the members of the management team believed they had no other choice.

The worker erupted at Mrs.~Polum, asking how he was supposed to support his family now.

Mrs.~Polum mumbled an apology, guilt washing over her. For employees at Ram, she said later,
``you're part of my family. We've been together. I'm supposed to protect you.''

She tries to remind herself that she is doing what is necessary to keep the company afloat. But
sometimes that is of little solace.

Hurrying back to her office that morning, Mrs.~Polum shrugged off a colleague who asked if she was
O.K. She closed the door and slumped into the corner in tears.

``I'm hoping we've found the bottom,'' she said. ``I feel like we can see the light. There is a
chance we're done. We've seen the worst. I'm hoping.''

\section{China Reports Student Toll for Quake}

BEIJING -- After a year of obfuscation, the authorities on Thursday released the first official
tally of student deaths from the earthquake last May, saying that 5,335 children either were dead or
remained missing. An additional 546 were left disabled, they said.

Previous estimates placed the number of students who died in the collapse of school buildings during
the May 12 earthquake in Sichuan Province as high as 10,000.

During a news conference in Chengdu, Tu Wentao, head of the Sichuan provincial education department,
said that the student death figures were accurate.

``These numbers were reached through legal methods,'' he said. ``We have wide agreement on these
numbers.'' The issue of student deaths remains a contentious one here. The parents of children who
perished in the rubble of classrooms say the buildings were poorly constructed; the government has
largely quashed the issue by harassing or detaining those who insist on pushing the matter.

With the first anniversary of the quake fast approaching, the government has stepped up its campaign
to silence those who have been calling for a full accounting of why so many schools failed while
adjacent structures remained undamaged. In recent days, several parents whose children died -- and
who have refused to stay quiet -- said they had been placed under heightened surveillance, and some
foreign journalists who have tried to interview grieving parents have been detained.

The newly released numbers did little to quell critics. Ai Weiwei, an artist who is one of China's
best-known gadflies, said the figures were ``meaningless'' because they lacked specifics, like
names, ages and places of death. Although Mr.~Ai, who says he has been compiling his own detailed
list of the dead, estimates the final figure at around 6,000, he said more precise information from
the government might shed some light on the issue of substandard school construction or at least
show that the government was truly willing to tackle that fundamental concern.

``The government is trying to escape from the accountability on this matter by postponing or
rejecting the publication of details,'' he said.

Government officials say that 68,712 people died during the earthquake. Another 17,921 are listed as
missing but are presumed to be dead. According to the official media, 7,000 classrooms and dormitory
rooms collapsed during the quake.

Over the past year, a number of other people who sought to draw attention to the questions over
school construction have been punished or are awaiting trial, including Liu Shakun, a teacher who
was charged with ``disseminating rumors and disrupting social order'' after he posted photographs of
destroyed schools on the Internet.

The parents of 126 children who died in Fuxin No.~2 Primary School in the town of Mianzhu have been
among the most outspoken -- and among those most watched by the government. Sang Jun, who lost his
11-year-old son last year, said in a telephone interview that since Monday, dozens of parents had
been watched by about 200 security officers and minor officials sent by the Mianzhu government.

Mr.~Sang said an official told him that the parents would be closely monitored until after the
anniversary of the quake, and that any contact with foreign journalists was considered ``unfavorable
to China.''

He and his wife met with a reporter and photographer from The New York Times last week to discuss
their grievances. At the time, several parents were willing to secretly meet with journalists in
their homes.

Mr.~Sang said a half-dozen people, mostly village officials, had been watching his home, but he
drove them away after an argument on Wednesday morning. His house, on the edge of a wheat field, is
close to the site of the Fuxin school, where new school buildings are being built with money from
Taiwan. Mr.~Sang said the local government was paying the minders \$22 per day to keep an eye on the
parents, the school site and roads leading into the village.

Foreign journalists trying to interview parents have been detained, and some had equipment broken by
security officers in recent weeks, according to a report issued Wednesday by the Foreign
Correspondents' Club of Beijing, which advocates greater press freedoms in China. A reporter for The
Financial Times was punched Tuesday by a thug, possibly a security officer in regular clothes, while
conducting interviews around Mianzhu, the group said. After he and his colleagues retreated to their
car, they were surrounded by a dozen hostile men, one of whom tried to hit a Chinese news assistant.

Liu Xiaoying, whose daughter died in the Fuxin school, said a French television crew was detained by
security officers and led out of the temporary housing camp where Ms.~Liu lives after the officers
learned that the crew was interviewing her.

On Monday a group of parents from Mianzhu secretly traveled to Beijing and filed a petition at the
central government's petitioning center. It was the third time parents from Mianzhu had tried to
file a complaint in the capital.

Employees at the petitioning center immediately notified officials in Mianzhu, who then had people
in Beijing detain the parents and escort them back to Sichuan, Mr.~Sang said. After the parents
returned to Mianzhu on Wednesday morning, they were immediately forced into a hospital and told they
would be checked for swine flu, Mr.~Sang said.

When asked about the reports of harassment, the head of the publicity department in Mianzhu, a woman
who gave her name as Ms.~Xu, said, ``As far as I know, this doesn't happen.''

\section{A Dividing Line Through the Financial Landscape}

At one bank in Alabama, the problem is a construction bust. At two in Ohio, the trouble is real
estate. And in New York, at the giant Citigroup, the worry is credit cards.

The stress tests released by the Obama administration Thursday painted a broad montage of the
troubles in the nation's banking industry and, for the first time, drew a stark, official line
through the new landscape of American finance.

On one side are institutions like JPMorgan Chase and Goldman Sachs, which regulators deemed stronger
than their peers -- perhaps strong enough to repay billions of bailout dollars and wriggle free of
government control.

On the other side are weaker institutions, among them Bank of America, which now confront the
daunting challenge of raising additional capital on their own or accepting increased government
ownership, along with whatever strings might be attached. Time is short: the banks have only until
June 8 to draw up their plans for regulators. As the results of the tests streamed in from the
Federal Reserve late Thursday, banks began racing to raise newly required capital.

Broadly speaking, the test results suggested that the banking industry is in better shape than many
had feared. Of the nation's 19 largest banks, which sit atop two-thirds of all deposits, regulators
gave nine a clean bill of health. The remaining 10 were ordered to raise a combined \$75 billion in
equity capital as a buffer against potential losses should the economy deteriorate. That amount is
far less than even the most optimistic forecast.

But the tests left some crucial questions unanswered. Once the banks satisfy regulators, will they
start making more of the loans that businesses and consumers need? And what happens if this
recession turns out to be even worse than regulators' worst-case assumptions -- and the banks'
losses start ballooning again?

``Everybody should breathe a sigh of relief,'' said Peter J.~Solomon, who runs a boutique investment
bank and early in his career worked at Lehman Brothers. ``Now they question is, So what? Will they
lend?''

The results shined an uncomfortable spotlight on the most troubled financial institutions: GMAC, the
finance arm of General Motors; Bank of America; Wells Fargo; and Citigroup. Several large regional
banks were also deemed to need large sums of capital.

Administration officials, as well as banking executives, hope the banks can plug many of the gaps by
selling stock to the public or selling businesses. If that is not enough, the banks could ask the
government to convert preferred shares purchased last autumn to common stock. Such a move would not
involve giving banks additional taxpayer money, although it would represent a riskier deal for the
government. It also raises the prospect that the government would become a substantial shareholder
in several large banks.

Wells Fargo announced ambitious plans to raise capital in hopes of avoiding increased government
ownership. Wells, which regulators determined needed \$13.7 billion, said it would seek to sell \$6
billion of stock. Morgan Stanley, which was ordered to raise \$1.8 billion in common equity, said it
would sell \$2 billion in new shares.

Bank of America, which faces a gap of roughly \$34 billion, said it would make a common stock
offering on Friday and would offer an exchange for some of its preferred stock later this month. It
said it expected to raise \$17 billion this way and also planned to sell assets worth \$10 billion.
The bank said it would fulfill the rest of \$34 billion needed through earnings in next few months.

Kenneth D.~Lewis, the bank's embattled chief executive, hailed the results of the stress test as a
sign of his bank's health.

``The bogey is large but we have significant opportunities,'' he said in a conference call Thursday
evening. ``We do not need new government money, and we do not intend to convert the existing TARP
money that we have. In fact our game plan is to get the government out of our bank as quickly as
possible.''

Mr.~Lewis noted that shareholders last week expressed concerns at the annual meeting and stripped
him of his chairman's title. ``As you might imagine, it was a humbling experience,'' he said.

``Frankly there's been a lot of noise around our company,'' Mr.~Lewis said. ``We hope these changes
will help us quiet that noise.''

Citigroup, which for many has come to symbolize the problems plaguing the financial industry, was
has already been moving quickly to address its problems. Regulators determined that the bank must
raise \$5.5 billion, on top of the converting \$45 billion in bailout funds to common stock, which
would give the government a 36 percent stake.

Vikram S.~Pandit, the bank's chief executive, said he would expand the company's offer to exchange
preferred shares of stock for common stock to a broader assortment of private investors. He also
said he would speed up plans to sell businesses.

``The government's stress test was a rigorous process that assessed our capital and confirms our
view that Citi's plans and actions will give it the financial strength to weather an adverse stress
scenario,'' Mr.~Pandit said in a statement.

Citigroup has already accepted the government as a key owner. In February, it announced that it
would convert just over half of the \$45 billion it has received from the government into common
stock, a step that will give taxpayers an ownership stake of 36 percent.

At the Regions Financial Corporation, the largest bank in Alabama, was ordered to increase its
common equity by \$2.5 billion. The bank is exposed to steep losses on construction loans. But
Regions questioned regulators' determination, saying it believed the projected losses were
``unrealistically high.'' Nonetheless, Regions said it would comply and that it was considering its
options.

SunTrust Banks, a large lender in Georgia that has also been hurt by the construction bust, was told
to raise \$2.2 billion. In Ohio, Keycorp and Fifth Third Bancorp, both of which expanded into
Florida at the height of the real estate bubble and must now raise additional capital to cover their
projected losses, said they were evaluating their options.

But a handful of stronger banks are pulling away from their weaker peers and emerging as
institutions that could dominate the industry.

Goldman Sachs, for instance, said Thursday that it hopes to return the \$10 billion it received from
the government as soon as possible, and regulators blessed Goldman's capital position as it is.
JPMorgan Chase, which was also deemed to have enough capital, is pushing to return bailout money as
well.

For Washington and Wall Street, the overarching question is whether investors and depositors will
take solace from these results. Banks and administration officials are eager to persuade private
investors that the banks are stable enough to invest in. While some analysts question the rigor of
the tests, Timothy F.~Geithner, the Treasury secretary, has insisted the exams will foster
confidence in the banks.

``The Fed is hoping the whole process will be explained to investors and there will be enough faith
in the legitimacy of the whole stress tests that they'll put to rest all of the speculation,'' said
Tanya Azarchs, the bank rating analyst at Standard \& Poor's. ``The golden ring here, if you can
catch it, is confidence.''

\section{G.M.'s Cash Reserves Are Dwindling}

DETROIT -- Even after receiving \$15.4 billion in federal loans, General Motors is once again on the
brink of financial collapse.

The automaker's first-quarter earnings released Thursday showed that G.M. was losing more money and
sales than it was in late December, when the government began its bailout.

With its cash reserves down to the bare minimum and its revenue plunging, G.M. seems more certain
each day to be heading toward a bankruptcy filing.

The company is entering the final stages of producing, by June 1, a restructuring plan that is
acceptable to President Obama and his auto task force.

But its financial problems continue to deepen. G.M. is spending more than \$113 million a day than
it is taking in from sales of its vehicles around the world.

Even if G.M. satisfies Mr.~Obama's other requirements -- including reaching a new cost-cutting deal
with its unions and persuading its bondholders to agree to sharply reduce debt -- the company looks
less viable, not more, than it did five months ago.

With Chrysler already in bankruptcy, many industry analysts say they think G.M. is next.

``It's looking like a real high probability,'' said Brett D.~Hoselton, an analyst with KeyBanc
Capital Markets. ``Chrysler is the best indicator at this point of where we're heading with G.M.''

Investors appear to be arriving at the same conclusion. G.M.'s stock dropped to \$1.60, down 3
percent, on Thursday, and its market capitalization fell below \$1 billion.

In the first quarter, G.M. lost \$5.9 billion, excluding special items, during the three-month
period that ended in March. Its revenue fell to \$22.4 billion -- a 47 percent drop compared with
the first quarter of 2008. It was the company's eighth consecutive quarterly loss.

G.M. said it built 1.33 million vehicles globally in the quarter, 903,000 fewer than it did in the
same period last year.

The company's chief financial officer, Ray Young, called the drop in production ``a staggering
number,'' and said consumers were showing increasing concern about G.M. products because of the
potential for bankruptcy.

If G.M. does have to file for bankruptcy protection, Mr.~Young emphasized Thursday that the
automaker would need to emerge quickly from court proceedings to keep sales from falling even more.

``Once you start losing revenues, you get yourself into a vicious cycle from which you cannot
recover,'' Mr.~Young said in a conference call with analysts and reporters. ``We prefer to
restructure outside of bankruptcy, but if we have to go in, we need to go in and out quickly.''

G.M. ended the quarter with \$11.6 billion in cash, roughly the level of liquidity it has said it
needs to keep operating and paying its bills.

In late April, the government lent G.M. an additional \$2 billion and committed to providing the
company with enough working capital to reach the June 1 deadline.

During the quarter, G.M. depleted its cash reserves by \$10.2 billion. That is twice the company's
rate of spending in the fourth quarter.

Industry analysts said the results would have been worse if the company hadn't already taken drastic
cost-cutting steps, including shedding 20,000 jobs since the middle of last year.

``It almost didn't matter what General Motors' first-quarter results looked like,'' said Shelly
Lombard, an analyst with the bond research firm Gimme Credit. ``With auto sales off dramatically,
the market expected G.M.'s results to be terrible and one good quarter wouldn't have been enough to
convince the market of a turnaround anyway.''

The quarterly results, once so closely watched for signs of progress at G.M., were viewed more as
evidence of its slide toward bankruptcy.

``The numbers are almost inconsequential,'' said Kip D.~Penniman Jr., a credit analyst with KDP
Investment Advisors. ``We know that G.M. is presently unsustainable.''

G.M.'s chief executive, Fritz Henderson, said Thursday that the company was at a ``defining moment''
in its 100-year history, and was continuing to work toward fulfilling the government's requirements
for future assistance.

The automaker this week reopened talks with the United Automobile Workers on using stock to finance
half of its \$20 billion in retiree health care obligations.

The U.A.W. reached a similar deal with Chrysler before it entered bankruptcy last week. But G.M.'s
plans to cut an additional 21,000 union jobs has drawn fire from union officials, and could make the
negotiations tougher than at Chrysler.

G.M. is also offering stock to the thousands of individual bondholders who collectively hold \$27
billion in the company's debt. G.M. has said it needs to cut the debt by at least \$24 billion.

Even Mr.~Henderson has questioned whether G.M. has a realistic chance to get 90 percent of its
bondholders to convert their debt to equity. Analysts discount the possibility almost entirely.

``G.M. will never get anywhere close to 90 percent of the bondholders,'' Mr.~Penniman said, adding
that a determining factor on a G.M. bankruptcy might be how quickly Chrysler's case moved forward.

\section{In Strategy Shift, G.E. Plans Lower-Cost Health Products}

General Electric is shifting the strategy in its \$17 billion-a-year health equipment and technology
business, seeking to broaden its reach with more lower-cost products.

The move is part of a wide-ranging marketing and business initiative announced on Thursday, called
``healthymagination.'' General Electric's health business, known for its medical imaging and
diagnostic machines, is struggling in a weak economy.

The hospitals and clinics that buy such sophisticated and costly equipment are reducing capital
spending in the downturn. In the first quarter, revenue for G.E.'s health care business fell 9
percent, to \$3.55 billion, and its operating profit fell 22 percent, to \$411 million.

``The high end in health care is never going to go away,'' Jeffrey R.~Immelt, G.E.'s chief
executive, said in an interview on Thursday. ``But this will make us broader in terms of price
points and offerings.''

As part of the new health strategy, G.E. said it would invest \$3 billion to develop at least 100
product and services innovations by 2015. The criteria for the new products, G.E. said, will be that
they lower cost, increase access and improve quality by 15 percent. The threshold measurements will
be done by Oxford Analytica, a consulting and analysis firm.

At a news conference in Washington, G.E. showed off some of the new products that it planned to
bring to market, including a hand-held ultrasound machine, decision-support software for physicians
developed with the Mayo Clinic and Intermountain Healthcare, and a bassinet with an overhead heater
for keeping newborns warm, a product only for poorer nations.

From \$200 million to \$300 million a year of the health research and development spending would be
additional money, Mr.~Immelt said, while the rest would be from the health unit's existing research
budget.

Broadly, the innovations are intended to address four needs: accelerating the adoption of health
information technology like computerized patient records, delivering high technology medical
equipment at lower prices, improving the access to health care by underserved communities and
supporting consumer-driven preventive care.

``This is going to force G.E. to wrestle with big issues -- technical ones and beyond,'' said Noel
M.~Tichy, a professor of management and organizations at the University of Michigan business school.
``It's a real broadening, and very different from selling MRI machines to big hospitals.''

The G.E. initiative will involve businesses beyond its health care division. In June, the company's
MSNBC cable television network plans to introduce a daily health program, anchored by the NBC News
chief medical editor, Dr.~Nancy L.~Snyderman.

G.E., which is based in Fairfield, Conn., also intends to improve the health of the 600,000 workers
and retirees it insures, using education, incentives and preventive health programs. One goal is to
keep the increase in G.E. health costs below the general inflation rate.

The health campaign is broadly modeled on the company's ``ecomagination'' initiative, begun in 2005.
It is focused on making G.E. and its products more environmentally friendly. Since 2005, the company
has invested \$1.5 billion on research and development for ``green'' products, tripled sales of such
products to \$18 billion and reduced the company's carbon footprint by 8 percent, saving \$100
million a year.

\section{Google’s Strength May Be Part of Microsoft Defense Strategy}

BERLIN -- Microsoft will argue against a European Commission proposal that it promote competing
browsers in its Windows operating system on the ground that such a move would strengthen its rival
Google's dominance in the global search-advertising market, according to a person with direct
knowledge of Microsoft's legal defense.

The company will make the argument at a June hearing in Brussels as part of an antitrust inquiry
about the packaging of its Internet Explorer browser with Windows, which powers more than 90 percent
of the world's personal computers.

The person with knowledge of Microsoft's legal strategy, who declined to be named because he was not
authorized to speak publicly about the case, said that Microsoft would outline what it saw as the
damaging effects to the search-advertising industry of incorporating competing browsers -- like
Firefox from Mozilla or Chrome from Google -- into Windows.

Mozilla derives the bulk of its income from the fees it receives for driving Web traffic to Google's
search engine. Google makes most of its revenue with advertising aimed at search results. Google
dominates the search advertising sector.

``Not only would Google's browser Chrome suddenly be on all Windows PCs, but it would strengthen
Google's dominance in search advertising,'' said this person.

Google, in a statement, did not directly respond to Microsoft's argument.

``We believe more competition will mean greater innovation on the Web and a better user experience
for people everywhere,'' said William Echikson, a Google spokesman in Brussels.

By aiming at its rival, Microsoft underscores the stakes for both it and Google, which are
increasingly clashing as personal computing moves from software to the Web. Aside from competing in
e-mail, online advertising and search, Google also makes an operating system used on mobile
computers and cellphones.

The new European complaint over browsers, like the previous nine-year antitrust case centering on
Microsoft's media player and computer-server coding, has forced the biggest names in the global
software and Web businesses to choose sides.

The complaint was brought in December 2007 by Opera, a tiny Norwegian software maker with 510
employees. Since then, Microsoft's main competitors -- Google, Mozilla, Sun Microsystems, Nokia,
I.B.M., Adobe, Red Hat and Corel -- have all signed on directly or through industry groups to argue
as co-complainants.

Google, Mozilla and the European Committee for Interoperable Systems, a group based in Brussels
representing the other Microsoft rivals, each plan to testify during the closed-door hearing next
month before Michael Albers, a hearing officer in the European Commission's competition section.

On Jan.~15, the commission informed Microsoft that it might order the software maker to preinstall
competitors' browsers in Windows, while simultaneously disabling Internet Explorer. Microsoft has
warned investors that the case could lead to a ``significant'' fine. Microsoft paid more than 1
billion euros in fines and penalties fighting the first antitrust case, which it abandoned in 2007
after losing an important legal appeal.

According to the person, Microsoft will argue that Internet browsing is inseparable from the Windows
operating system. Microsoft will also emphasize that consumers can download and use any competing
browser with Windows, and that Internet Explorer's share of the browser market has been falling
steadily.

An estimated 67 percent of computer users around the world use Microsoft's Internet Explorer,
according to Net Applications, a researcher in Aliso Viejo, Calif. In Europe, the Firefox 3 browser
has overtaken Internet Explorer 7 as the most-used browser, according to the Web analytics firm
StatCounter. But it still trails Microsoft by 10 percentage points when all versions of the Explorer
are counted.

The interoperable systems committee has said that it intended to present data at the hearing showing
that Microsoft's browser market share was closer to 85 percent in Europe.

\section{No Slowing in Cash Flow for `Idol'}

LOS ANGELES -- Can a television show with a shrinking audience actually increase its revenue?

In the case of ``American Idol,'' the answer is yes, and the manner that the show's creators are
going about it could change the way that producers and broadcast networks think about the longevity
of reality television programs.

Despite losing viewers in each of the last three years, ``American Idol,'' still the top-rated
series on television, is generating ever-growing profits for its creators and, it appears, for Fox
Broadcasting through brand extensions, marketing arrangements and licensing fees.

The deals, which include products as disparate as ice cream and trading cards, as well as the more
familiar partnerships with iTunes and AT\&T, have driven tremendous growth in the profitability of
``American Idol,'' according to the public financial statements of the parent of 19 Entertainment,
the company founded by Simon Fuller, the creator of the show.

Those financial statements include documents that break down how revenue is split among Fox, 19
Entertainment and FremantleMedia, the production company behind the show.

``They are just raking in the money, hand over fist,'' said Brad Adgate, a senior vice president for
research at Horizon Media. ``There are certain things that 'American Idol' does not do well, like
selling DVDs. But the whole industry is moving toward a different model where deciding whether
something is a hit or not is not just based on how much you charge for 30 seconds of advertising.''

While everyone in television, from Ben Silverman, the co-chairman of NBC, to the youth-centric
programmers at MTV, is trying to find ways to increase profits as television audiences decline,
``American Idol'' is taking its cues not from the traditional broadcasting playbook but from the
National Football League.

``We have learned the lessons of the sports leagues in that they have all these ancillary revenue
streams,'' said Robert F.~X. Sillerman, chief executive of CKX Inc., the parent of 19 Entertainment.
``And frankly, we're just beginning.''

David Lunar, a senior vice president for interactive and consumer products at FremantleMedia, which
produces the program in partnership with 19 Entertainment, said this year that the company had
experienced ``eight straight years of growth'' in licensing and merchandising.

That is particularly noteworthy given that, according to Nielsen Media Research, the average
audience for ``American Idol'' peaked three years ago, at more than 30 million viewers an episode.
This season, each episode is attracting about 25 million viewers an episode.

Advertising revenue, which primarily benefits Fox Broadcasting, has grown in each of the last three
years, according to TNS Media Intelligence, to \$903 million last year. That is nearly double the
level of three years earlier. Those figures do not include Fox's ancillary sponsorship deals and
other income, like royalties it receives from the sale of music performances by ``Idol''
contestants.

Even though ``American Idol's'' ratings have declined, Fox has put more hours of the show on the air
and has been able to charge a higher rate as its ratings lead over other shows has grown.

The revenue and profits of 19 Entertainment have grown even more quickly, according to the financial
statements of CKX. Its revenues from ``American Idol'' alone grew to \$96 million last year from
\$67 million two years earlier, with gross profit margins expanding to 77 percent from 69 percent in
that span.

Including revenue from ``Idol'' programming in other countries, from music sales related to all the
``Idol'' shows and from ``So You Think You Can Dance,'' also on Fox, 19 Entertainment produced
revenue of \$223 million last year, up from \$151 million two years earlier.

Fremantle does not break out its revenue, and a spokesman for the company declined to comment on its
financial performance. But contracts included in the CKX financial statements stipulate that
Fremantle receive one-third to one-half of the various revenue streams that benefit 19
Entertainment.

Even though the overall audience is declining, ``American Idol'' appears to be in little danger of
losing its crown as the top-rated television series. In the 2003-4 season, the first in which
``Idol'' was the top-rated prime-time series, its lead over the second place show was about 7
percent. That margin has grown every year since and this year is 66 percent.

Mike Darnell, the president of alternative entertainment for Fox, who oversees ``Idol'' and other
reality shows for the network, said that ``Idol'' could lose 12 percent of its audience every season
and still be among the top 10 shows on television in 2016 -- even if every other show on television
maintained all of its current audience.

Fox, too, benefits from the brand extensions. Its Fox Reality Channel televises ``American Idol
Rewind,'' which repeats segments from earlier seasons along with previously unseen footage of
noteworthy contestants, and it has turned drama series like ``House'' and ``Fringe'' into hits by
pairing them with ``Idol'' for a season or two.

While some programs, like ``The Simpsons,'' generate large revenue from merchandising and marketing
deals separate from the broadcast, ``they don't have the audience delivery or all the hoopla of
'American Idol,' '' said Mr.~Adgate. ``I don't think there are any reality shows that have a
similar, off-network shelf life.''

Mr.~Darnell declined to comment on just how much Fox benefited from those and similar efforts. ``But
I can say the brand extension is huge for this company and for all the companies involved with the
show,'' he said.

There is evidence that the efforts, if not drawing more people to the program, are certainly getting
viewers more engaged. Last week, viewers cast 64 million votes, the most ever for a nonfinale
episode.

Mr.~Sillerman cited that jump as evidence that ``American Idol'' could run indefinitely.

``If you only had a 10-year-old, you would probably say the impact of 'Sesame Street' has
diminished, because it's not on in your house anymore,'' Mr.~Sillerman said. ``But they're making
new 5-year-olds every year.''

\section{China Outpaces U.S. in Cleaner Coal-Fired Plants}

TIANJIN, China -- China's frenetic construction of coal-fired power plants has raised worries around
the world about the effect on climate change. China now uses more coal than the United States,
Europe and Japan combined, making it the world's largest emitter of gases that are warming the
planet.

But largely missing in the hand-wringing is this: China has emerged in the past two years as the
world's leading builder of more efficient, less polluting coal power plants, mastering the
technology and driving down the cost.

While the United States is still debating whether to build a more efficient kind of coal-fired power
plant that uses extremely hot steam, China has begun building such plants at a rate of one a month.

Construction has stalled in the United States on a new generation of low-pollution power plants that
turn coal into a gas before burning it, although Energy Secretary Steven Chu said Thursday that the
Obama administration might revive one power plant of this type. But China has already approved
equipment purchases for just such a power plant, to be assembled soon in a muddy field here in
Tianjin.

``The steps they've taken are probably as fast and as serious as anywhere in power-generation
history,'' said Hal Harvey, president of ClimateWorks, a group in San Francisco that helps finance
projects to limit global warming.

Western countries continue to rely heavily on coal-fired power plants built decades ago with
outdated, inefficient technology that burn a lot of coal and emit considerable amounts of carbon
dioxide. China has begun requiring power companies to retire an older, more polluting power plant
for each new one they build.

Cao Peixi, the president of the China Huaneng Group, the country's biggest state-owned electric
utility and the majority partner in the joint venture building the Tianjin plant, said his company
was committed to the project even though it would cost more than conventional plants.

``We shouldn't look at this project from a purely financial perspective,'' he said. ``It represents
the future.''

Without doubt, China's coal-fired power sector still has many problems, and global warming gases
from the country are expected to continue increasing. China's aim is to use the newest technologies
to limit the rate of increase.

Only half the country's coal-fired power plants have the emissions control equipment to remove
sulfur compounds that cause acid rain, and even power plants with that technology do not always use
it. China has not begun regulating some of the emissions that lead to heavy smog in big cities.

Even among China's newly built plants, not all are modern. Only about 60 percent of the new plants
are being built using newer technology that is highly efficient, but more expensive.

With greater efficiency, a power plant burns less coal and emits less carbon dioxide for each unit
of electricity it generates. Experts say the least efficient plants in China today convert 27 to 36
percent of the energy in coal into electricity. The most efficient plants achieve an efficiency as
high as 44 percent, meaning they can cut global warming emissions by more than a third compared with
the weakest plants.

In the United States, the most efficient plants achieve around 40 percent efficiency, because they
do not use the highest steam temperatures being adopted in China. The average efficiency of American
coal-fired plants is still higher than the average efficiency of Chinese power plants, because China
built so many inefficient plants over the past decade. But China is rapidly closing the gap by using
some of the world's most advanced designs.

After relying until recently on older technology, ``China has since become the major world market
for advanced coal-fired power plants with high-specification emission control systems,'' the
International Energy Agency said in a report on April 20.

China's improvements are starting to have an effect on climate models. In its latest annual report
last November, the I.E.A. cut its forecast of the annual increase in Chinese emissions of global
warming gases, to 3 percent from 3.2 percent, in response to technological gains, particularly in
the coal sector, even as the agency raised slightly its forecast for Chinese economic growth. ``It's
definitely changing the baseline, and that's being taken into account,'' said Jonathan Sinton, a
China specialist at the energy agency.

But by continuing to rely heavily on coal, which supplies 80 percent of its electricity, China
ensures that it will keep emitting a lot of carbon dioxide; even an efficient coal-fired power plant
emits twice the carbon dioxide of a natural gas-fired plant.

Perhaps the biggest question now is how much further China can go beyond the recent steps. In
particular, how fast will it move toward power plants that capture their emissions and store them
underground or under the seafloor?

That technology could, in theory, create power plants that contribute virtually nothing to global
warming. Many countries hope to develop such plants, though progress has been halting; Energy
Secretary Chu has promised steps to speed up the technology in the United States.

China has just built a small, experimental facility near Beijing to remove carbon dioxide from power
station emissions and use it to provide carbonation for beverages, and the government has a short
list of possible locations for a large experiment to capture and store carbon dioxide. But so far,
it has no plans to make this a national policy.

China is making other efforts to reduce its global warming emissions. It has doubled its total wind
energy capacity in each of the past four years, and is poised to pass the United States as soon as
this year as the world's largest market for wind power equipment. China is building considerably
more nuclear power plants than the rest of the world combined, and these do not emit carbon dioxide
after they are built.

But coal remains the cheapest energy source in China by a wide margin. China has the world's
third-largest coal reserves, after the United States and Russia.

``No matter how much renewable or nuclear is in the mix, coal will remain the dominant power
source,'' said Ashok Bhargava, a China energy expert at the Asian Development Bank in Manila.

Another problem is that China has finally developed the ability to build high-technology power
plants only at the end of a national binge of building lower-tech coal-fired plants. Construction is
now slowing because of the economic slump.

By adopting ``ultra-supercritical'' technology, which uses extremely hot steam to achieve the highest
efficiency, and by building many identical power plants at the same time, China has cut costs
dramatically through economies of scale. It now can cost a third less to build an
ultra-supercritical power plant in China than to build a less efficient coal-fired plant in the
United States.


\section{CYBERWAR:U.S. Steps Up Effort on Digital Defenses}

This article was reported by David E.~Sanger, John Markoff and Thom Shanker and written by
Mr.~Sanger.

When American forces in Iraq wanted to lure members of Al Qaeda into a trap, they hacked into one of
the group's computers and altered information that drove them into American gun sights.

When President George W.~Bush ordered new ways to slow Iran's progress toward a nuclear bomb last
year, he approved a plan for an experimental covert program -- its results still unclear -- to bore
into their computers and undermine the project.

And the Pentagon has commissioned military contractors to develop a highly classified replica of the
Internet of the future. The goal is to simulate what it would take for adversaries to shut down the
country's power stations, telecommunications and aviation systems, or freeze the financial markets
-- in an effort to build better defenses against such attacks, as well as a new generation of online
weapons.

Just as the invention of the atomic bomb changed warfare and deterrence 64 years ago, a new
international race has begun to develop cyberweapons and systems to protect against them.

Thousands of daily attacks on federal and private computer systems in the United States -- many from
China and Russia, some malicious and some testing chinks in the patchwork of American firewalls --
have prompted the Obama administration to review American strategy.

President Obama is expected to propose a far larger defensive effort in coming days, including an
expansion of the \$17 billion, five-year program that Congress approved last year, the appointment
of a White House official to coordinate the effort, and an end to a running bureaucratic battle over
who is responsible for defending against cyberattacks.

But Mr.~Obama is expected to say little or nothing about the nation's offensive capabilities, on
which the military and the nation's intelligence agencies have been spending billions. In interviews
over the past several months, a range of military and intelligence officials, as well as outside
experts, have described a huge increase in the sophistication of American cyberwarfare capabilities.

Because so many aspects of the American effort to develop cyberweapons and define their proper use
remain classified, many of those officials declined to speak on the record. The White House declined
several requests for interviews or to say whether Mr.~Obama as a matter of policy supports or
opposes the use of American cyberweapons.

The most exotic innovations under consideration would enable a Pentagon programmer to
surreptitiously enter a computer server in Russia or China, for example, and destroy a ``botnet'' --
a potentially destructive program that commandeers infected machines into a vast network that can be
clandestinely controlled -- before it could be unleashed in the United States.

Or American intelligence agencies could activate malicious code that is secretly embedded on
computer chips when they are manufactured, enabling the United States to take command of an enemy's
computers by remote control over the Internet. That, of course, is exactly the kind of attack
officials fear could be launched on American targets, often through Chinese-made chips or computer
servers.

So far, however, there are no broad authorizations for American forces to engage in cyberwar. The
invasion of the Qaeda computer in Iraq several years ago and the covert activity in Iran were each
individually authorized by Mr.~Bush. When he issued a set of classified presidential orders in
January 2008 to organize and improve America's online defenses, the administration could not agree
on how to write the authorization.

A principal architect of that order said the issue had been passed on to the next president, in part
because of the complexities of cyberwar operations that, by necessity, would most likely be
conducted on both domestic and foreign Internet sites. After the controversy surrounding domestic
spying, Mr.~Bush's aides concluded, the Bush White House did not have the credibility or the
political capital to deal with the subject.

Electronic Vulnerabilities

Cyberwar would not be as lethal as atomic war, of course, nor as visibly dramatic. But when Mike
McConnell, the former director of national intelligence, briefed Mr.~Bush on the threat in May 2007,
he argued that if a single large American bank were successfully attacked ``it would have an
order-of-magnitude greater impact on the global economy'' than the Sept.~11, 2001, attacks.
Mr.~McConnell, who left office three months ago, warned last year that ``the ability to threaten the
U.S.~money supply is the equivalent of today's nuclear weapon.''

The scenarios developed last year for the incoming president by Mr.~McConnell and his coordinator
for cybersecurity, Melissa Hathaway, went further. They described vulnerabilities including an
attack on Wall Street and one intended to bring down the nation's electric power grid. Most were
extrapolations of attacks already tried.

Today, Ms.~Hathaway is the primary author of White House cyberstrategy and has been traveling the
country talking in vague terms about recent, increasingly bold attacks on the computer networks that
keep the country running. Government officials will not discuss the details of a recent attack on
the air transportation network, other than to say the attack never directly affected air traffic
control systems.

Still, the specter of an attack that could blind air traffic controllers and, perhaps, the
military's aerospace defense networks haunts military and intelligence officials. (The saving grace
of the air traffic control system, officials say, is that it is so old that it is not directly
connected to the Internet.)

Studies, with code names like Dark Angel, have focused on whether cellphone towers,
emergency-service communications and hospital systems could be brought down, to sow chaos.

But the theoretical has, at times, become real.

``We have seen Chinese network operations inside certain of our electricity grids,'' said Joel
F.~Brenner, who oversees counterintelligence operations for Dennis Blair, Mr.~McConnell's successor
as national intelligence director, speaking at the University of Texas at Austin this month. ``Do I
worry about those grids, and about air traffic control systems, water supply systems, and so on? You
bet I do.''

But the broader question -- one the administration so far declines to discuss -- is whether the best
defense against cyberattack is the development of a robust capability to wage cyberwar.

As Mr.~Obama's team quickly discovered, the Pentagon and the intelligence agencies both concluded in
Mr.~Bush's last years in office that it would not be enough to simply build higher firewalls and
better virus detectors or to restrict access to the federal government's own computers.

``The fortress model simply will not work for cyber,'' said one senior military officer who has been
deeply engaged in the debate for several years. ``Someone will always get in.''

That thinking has led to a debate over whether lessons learned in the nuclear age -- from the days
of ``mutually assured destruction'' -- apply to cyberwar.

But in cyberwar, it is hard to know where to strike back, or even who the attacker might be. Others
have argued for borrowing a page from Mr.~Bush's pre-emption doctrine by going into foreign
computers to destroy malicious software before it is unleashed into the world's digital bloodstream.
But that could amount to an act of war, and many argue it is a losing game, because the United
States is more dependent on a constantly running Internet system than many of its potential
adversaries, and therefore could suffer more damage in a counterattack.

In a report scheduled to be released Wednesday, the National Research Council will argue that
although an offensive cybercapability is an important asset for the United States, the nation is
lacking a clear strategy, and secrecy surrounding preparations has hindered national debate,
according to several people familiar with the report.

The advent of Internet attacks -- especially those suspected of being directed by nations, not
hackers -- has given rise to a new term inside the Pentagon and the National Security Agency:
``hybrid warfare.''

It describes a conflict in which attacks through the Internet can be launched as a warning shot --
or to pave the way for a traditional attack.

Early hints of this new kind of warfare emerged in the confrontation between Russia and Estonia in
April 2007. Clandestine groups -- it was never determined if they had links to the Russian
government -- commandeered computers around the globe and directed a fire hose of data at Estonia's
banking system and its government Web sites.

The computer screens of Estonians trying to do business with the government online were frozen, if
they got anything at all. It was annoying, but by the standards of cyberwar, it was child's play.

In August 2008, when Russia invaded Georgia, the cyberattacks grew more widespread. Georgians were
denied online access to news, cash and air tickets. The Georgian government had to move its Internet
activity to servers in Ukraine when its own servers locked up, but the attacks did no permanent
damage.

Every few months, it seems, some agency, research group or military contractor runs a war game to
assess the United States' vulnerability. Senior intelligence officials were shocked to discover how
easy it was to permanently disable a large power generator. That prompted further studies to
determine if attackers could take down a series of generators, bringing whole parts of the country
to a halt.

Another war game that the Department of Homeland Security sponsored in March 2008, called Cyber
Storm II, envisioned a far larger, coordinated attack against the United States, Britain, Canada,
Australia and New Zealand. It studied a disruption of chemical plants, rail lines, oil and gas
pipelines and private computer networks. That study and others like it concluded that when attacks
go global, the potential economic repercussions increase exponentially.

To prove the point, Mr.~McConnell, then the director of national intelligence, spent much of last
summer urging senior government officials to examine the Treasury Department's scramble to contain
the effects of the collapse of Bear Stearns. Markets froze, he said, because ``what backs up that
money is confidence -- an accounting system that is reconcilable.'' He began studies of what would
happen if the system that clears market trades froze.

``We were halfway through the study,'' one senior intelligence official said last month, ``and the
markets froze of their own accord. And we looked at each other and said, 'Our market collapse has
just given every cyberwarrior out there a playbook.' ''

Just before Mr.~Obama was elected, the Center for Strategic and International Studies, a policy
research group in Washington, warned in a report that ``America's failure to protect cyberspace is
one of the most urgent national security problems facing the new administration.''

What alarmed the panel was not the capabilities of individual hackers but of nations -- China and
Russia among them -- that experts believe are putting huge resources into the development of
cyberweapons. A research company called Team Cymru recently examined ``scans'' that came across the
Internet seeking ways to get inside industrial control systems, and discovered more than 90 percent
of them came from computers in China.

Scanning alone does no damage, but it could be the prelude to an attack that scrambles databases or
seeks to control computers. But Team Cymru ran into a brick wall as soon as it tried to trace who,
exactly, was probing these industrial systems. It could not determine whether military
organizations, intelligence agencies, terrorist groups, criminals or inventive teenagers were behind
the efforts.

The good news, some government officials argue, is that the Chinese are deterred from doing real
damage: Because they hold more than a trillion dollars in United States government debt, they have
little interest in freezing up a system they depend on for their own investments.

Then again, some of the scans seemed to originate from 14 other countries, including Taiwan, Russia
and, of course, the United States.

Bikini Atoll for an Online Age

Because ``cyberwar'' contains the word ``war,'' the Pentagon has argued that it should be the locus
of American defensive and offensive strategy -- and it is creating the kind of infrastructure that
was built around nuclear weapons in the 1940s and '50s.

Defense Secretary Robert M.~Gates is considering proposals to create a Cyber Command -- initially as
a new headquarters within the Strategic Command, which controls the American nuclear arsenal and
assets in space. Right now, the responsibility for computer network security is part of Strategic
Command, and military officials there estimate that over the past six months, the government has
spent \$100 million responding to probes and attacks on military systems. Air Force officials
confirm that a large network of computers at Maxwell Air Force Base in Alabama was temporarily taken
off-line within the past eight months when it was put at risk of widespread infection from computer
viruses.

But Mr.~Gates has concluded that the military's cyberwarfare effort requires a sharper focus -- and
thus a specific command. It would build the defenses for military computers and communications
systems and -- the part the Pentagon is reluctant to discuss -- develop and deploy cyberweapons.

In fact, that effort is already under way -- it is part of what the National Cyber Range is all
about. The range is a replica of the Internet of the future, and it is being built to be attacked.
Competing teams of contractors -- including BAE Systems, the Applied Physics Laboratory at Johns
Hopkins University and Sparta Inc.~-- are vying to build the Pentagon a system it can use to
simulate attacks. The National Security Agency already has a smaller version of a similar system, in
Millersville, Md.

In short, the Cyber Range is to the digital age what the Bikini Atoll -- the islands the Army
vaporized in the 1950s to measure the power of the hydrogen bomb -- was to the nuclear age. But once
the tests at Bikini Atoll demonstrated to the world the awesome destructive power of the bomb, it
became evident to the United States and the Soviet Union -- and other nuclear powers -- that the
risks of a nuclear exchange were simply too high. In the case of cyberattacks, where the results can
vary from the annoying to the devastating, there are no such rules.

The Deterrence Conundrum

During the cold war, if a strategic missile had been fired at the United States, screens deep in a
mountain in Colorado would have lighted up and American commanders would have some time to decide
whether to launch a counterattack. Today, when Pentagon computers are subjected to a barrage, the
origin is often a mystery. Absent certainty about the source, it is almost impossible to mount a
counterattack.

In the rare case where the preparations for an attack are detected in a foreign computer system,
there is continuing debate about whether to embrace the concept of pre-emption, with all of its
Bush-era connotations. The questions range from whether an online attack should be mounted on that
system to, in an extreme case, blowing those computers up.

Some officials argue that if the United States engaged in such pre-emption -- and demonstrated that
it was watching the development of hostile cyberweapons -- it could begin to deter some attacks.
Others believe it will only justify pre-emptive attacks on the United States. ``Russia and China
have lots of nationalistic hackers,'' one senior military officer said. ``They seem very, very
willing to take action on their own.''

Senior Pentagon and military officials also express deep concern that the laws and understanding of
armed conflict have not kept current with the challenges of offensive cyberwarfare.

Over the decades, a number of limits on action have been accepted -- if not always practiced. One is
the prohibition against assassinating government leaders. Another is avoiding attacks aimed at
civilians. Yet in the cyberworld, where the most vulnerable targets are civilian, there are no such
rules or understandings. If a military base is attacked, would it be a proportional, legitimate
response to bring down the attacker's power grid if that would also shut down its hospital systems,
its air traffic control system or its banking system?

``We don't have that for cyber yet,'' one senior Defense Department official said, ``and that's a
little bit dangerous.''

\section{CYBERWAR:Iranians and Others Outwit Net Censors}

The Iranian government, more than almost any other, censors what citizens can read online, using
elaborate technology to block millions of Web sites offering news, commentary, videos, music and,
until recently, Facebook and YouTube. Search for ``women'' in Persian and you're told, ``Dear
Subscriber, access to this site is not possible.''

Last July, on popular sites that offer free downloads of various software, an escape hatch appeared.
The computer program allowed Iranian Internet users to evade government censorship.

College students discovered the key first, then spread it through e-mail messages and file-sharing.
By late autumn more than 400,000 Iranians were surfing the uncensored Web.

The software was created not by Iranians, but by Chinese computer experts volunteering for the Falun
Gong, a spiritual movement that has beem suppressed by the Chinese government since 1999. They
maintain a series of computers in data centers around the world to route Web users' requests around
censors' firewalls.

The Internet is no longer just an essential channel for commerce, entertainment and information. It
has also become a stage for state control -- and rebellion against it. Computers are becoming more
crucial in global conflicts, not only in spying and military action, but also in determining what
information reaches people around the globe.

More than 20 countries now use increasingly sophisticated blocking and filtering systems for
Internet content, according to Reporters Without Borders, a Paris-based group that encourages
freedom of the press.

Although the most aggressive filtering systems have been erected by authoritarian governments like
those in Iran, China, Pakistan, Saudi Arabia and Syria, some Western democracies are also beginning
to filter some content, including child pornography and other sexually oriented material.

In response, a disparate alliance of political and religious activists, civil libertarians, Internet
entrepreneurs, diplomats and even military officers and intelligence agents are now challenging
growing Internet censorship.

The creators of the software seized upon by Iranians are members of the Global Internet Freedom
Consortium, based largely in the United States and closely affiliated with Falun Gong. The
consortium is one of many small groups developing systems to make it possible for anyone to reach
the open Internet. It is the modern equivalent of efforts by organizations like the Voice of America
to reach the citizens of closed countries.

Separately, the Tor Project, a nonprofit group of anticensorship activists, freely offers software
that can be used to send messages secretly or to reach blocked Web sites. Its software, first
developed at the United States Naval Research Laboratories, is now used by more than 300,000 people
globally, from the police to criminals, as well as diplomats and spies.

Political scientists at the University of Toronto have built yet another system, called Psiphon,
that allows anyone to evade national Internet firewalls using only a Web browser. Sensing a business
opportunity, they have created a company to profit by making it possible for media companies to
deliver digital content to Web users behind national firewalls.

The danger in this quiet electronic war is driven home by a stark warning on the group's Web site:
``Bypassing censorship may violate law. Serious thought should be given to the risks involved and
potential consequences.''

In this cat-and-mouse game, the cat is fighting back. The Chinese system, which opponents call the
Great Firewall of China, is built in part with Western technologies. A study published in February
by Rebecca MacKinnon, who teaches journalism at the University of Hong Kong, determined that much
blog censorship is performed not by the government but by private Internet service providers,
including companies like Yahoo China, Microsoft and MySpace. One-third to more than half of all
postings made to three Chinese Internet service providers were not published or were censored, she
reported.

When the Falun Gong tried to support its service with advertising several years ago, American
companies backed out under pressure from the Chinese government, members said.

In addition, the Chinese government now employs more than 40,000 people as censors at dozens of
regional centers, and hundreds of thousands of students are paid to flood the Internet with
government messages and crowd out dissenters.

This is not to say that China blocks access to most Internet sites; most of the material on the
global Internet is available to Chinese without censorship. The government's censors mostly censor
groups deemed to be state enemies, like the Falun Gong, making it harder for them to reach potential
members.

Blocking such groups has become more insidious as Internet filtering technology has grown more
sophisticated. As with George Orwell's ``Newspeak,'' the language in ``1984'' that got smaller each
year, governments can block particular words or phrases without users realizing their Internet
searches are being censored.

Those who back the ragtag opponents of censorship criticize the government-run systems as the
digital equivalent of the Berlin Wall.

They also see the anticensorship efforts as a powerful political lever. ``What is our leverage
toward a country like Iran? Very little,'' said Michael Horowitz, a fellow at the Hudson Institute
who advises the Global Internet Freedom Consortium. ``Suppose we have the capacity to make it
possible for the president of the United States at will to communicate with hundreds of thousands of
Iranians at no risk or limited risk? It just changes the world.''

The United States government and the Voice of America have financed some circumvention technology
efforts. But until now the Falun Gong has devoted the most resources, experts said, erecting a
system that allows the largest number of Internet users open, uncensored access.

Each week, Chinese Internet users receive 10 million e-mail messages and 70 million instant messages
from the consortium. But unlike spam that takes you to Nigerian banking scams or offers deals on
drugs like Viagra, these messages offer software to bypass the elaborate government system that
blocks access to the Web sites of opposition groups like the Falun Gong.

Shiyu Zhou, a computer scientist, is a founder of the Falun Gong's consortium. His cyber-war with
China began in Tiananmen Square in 1989. A college student and the son of a former general in the
intelligence section of the People's Liberation Army, he said he first understood the power of
government-controlled media when overnight the nation's student protesters were transformed from
heroes to killers.

``I was so disappointed,'' he said. ``People believed the government, they didn't believe us.''

He decided to leave China and study computer science in graduate school in the United States. In the
late 1990s he turned to the study of Falun Gong and then joined with a small group of technically
sophisticated members of the spiritual group intent on transmitting millions of e-mail messages to
Chinese.

Both he and Peter Yuan Li, another early consortium volunteer, had attended Tsinghua University --
China's Massachusetts Institute of Technology. Mr.~Li, the son of farmers, also came to the United
States to study computer science, then joined Bell Laboratories before becoming a full-time
volunteer.

The risks of building circumvention tools became clear in April 2006 when, Mr.~Li later told law
enforcement officials, four men invaded his home in suburban Atlanta, covered his head, beat him,
searched his files and stole two laptop computers. The F.B.I. has made no arrests in the case and
declined to comment. But Mr.~Li thinks China sent the invaders.

Early on, the group of dissidents here had some financial backing from the International
Broadcasting Bureau of the Voice of America for sending e-mail messages, but the group insists that
most of its effort has been based on volunteer labor and contributions.

The consortium's circumvention system works this way: Government censorship systems like the Great
Firewall can block access to certain Internet Protocol addresses. The equivalent of phone numbers,
these addresses are quartets of numbers like 209.85.171.100 that identify a Web site, in this case,
google.com. By clicking on a link provided in the consortium's e-mail message, someone in China or
Iran trying to reach a forbidden Web site can download software that connects to a computer abroad
that then redirects the request to the site's forbidden address.

The technique works like a basketball bank shot -- with the remote computer as the backboard and the
desired Web site as the basket. But government systems hunt for and then shut off such alternative
routes using a variety of increasingly sophisticated techniques. So the software keeps changing the
Internet address of the remote computer -- more than once a second. By the time the censors identify
an address, the system has already changed it.

China acknowledges that it monitors content on the Internet, but claims to have an agenda much like
that of any other country: policing for harmful material, pornography, treasonous propaganda,
criminal activity, fraud. The government says Falun Gong is a dangerous cult that has ruined the
lives of thousands of people.

Hoping to step up its circumvention efforts, the Falun Gong last year organized extensive lobbying
in Congress, which approved \$15 million for circumvention services.

But the money was awarded not to the Falun Gong consortium but to Internews, an international
organization that supports local media groups.

This year, a broader coalition is organizing to push for more Congressional financing of
anti-filtering efforts. Negotiations are under way to bring together dissidents of Vietnam, Iran,
the Uighur minority of China, Tibet, Myanmar, Cuba, Cambodia, Laos, as well as the Falun Gong, to
lobby Congress for the financing.

Mr.~Horowitz argues that \$25 million could expand peak usage to as many as 45 million daily
Internet users, allowing the systems to reach as many as 10 percent of the Web users in both China
and Iran.

Mr.~Zhou says his group's financing is money well spent. ``The entire battle over the Internet has
boiled down to a battle over resources,'' he said. ``For every dollar we spend, China has to spend a
hundred, maybe hundreds of dollars.''

As for the Falun Gong software, it proved a little too popular among Iranians. By the end of last
year the consortium's computers were overwhelmed. On Jan.~1, the consortium had to do some blocking
of its own: It shut down the service for all countries except China.

\section{CYBERWAR:Cadets Trade the Trenches for Firewalls}

WEST POINT, N.Y. -- The Army forces were under attack. Communications were down, and the chain of
command was broken.

Pacing a makeshift bunker whose entrance was camouflaged with netting, the young man in battle
fatigues barked at his comrades: ``They are flooding the e-mail server. Block it. I'll take the heat
for it.''

These are the war games at West Point, at least last month, when a team of cadets spent four days
struggling around the clock to establish a computer network and keep it operating while hackers from
the National Security Agency in Maryland tried to infiltrate it with methods that an enemy might
use. The N.S.A. made the cadets' task more difficult by planting viruses on some of the equipment,
just as real-world hackers have done on millions of computers around the world.

The competition was a final exam of sorts for a senior elective class. The cadets, who were computer
science and information technology majors, competed against teams from the Navy, Air Force, Coast
Guard and Merchant Marine as well as the Naval Postgraduate Academy and the Air Force Institute of
Technology. Each team was judged on how well it subdued the threats from the N.S.A.

The cyberwar games at West Point are just one example of a heightened awareness across the military
that it must treat the threat of a computer attack as seriously as it does an attack carried out by
a bomber or combat brigade. There is hardly an American military unit or headquarters that has not
been ordered to analyze the risk of cyberattacks to its mission -- and to train to counter them. If
the hackers were to succeed, they could change information on the network and cripple Internet
communications.

In the desert outside Las Vegas, in a series of inconspicuous trailers, some of the most highly
motivated hackers in the United States spend their days and nights probing the military's vast
computer networks for weaknesses to exploit.

These hackers -- many of whom got their start as teenagers devoted to computer screens in their
basements -- have access to the latest in attack software. Some of it was developed by cryptologists
at the N.S.A., the nation's largest intelligence agency, where most of the government's talent for
breaking and making computer codes resides.

The hackers have an official name -- the 57th Information Aggressor Squadron -- and a real home,
Nellis Air Force Base.

The Army last year created its own destination for computer experts, the Network Warfare Battalion,
where many of the cadets in the cyberwar games hope to be assigned. But even so, the ranks are still
small.

The Defense Department today graduates only 80 students a year from its cyberwar schools, causing
Defense Secretary Robert M.~Gates to complain that the Pentagon is ``desperately short of people who
have capabilities in this area in all the services, and we have to address it.'' Under current
Pentagon budget proposals, the number of students cycled through the schools will be quadrupled in
the next two years.

Part of the Pentagon's effort to increase the military's capabilities are the annual cyberwar games
played at the nation's military academies, including West Point, where young cadets in combat boots
and buzz cuts talk megabytes instead of megatons on a campus dotted with statues of generals,
historic armaments and old stone buildings.

While the Pentagon has embraced the need for offensive cyberwarfare, there were no offensive
maneuvers in the games last month, said Col.~Joe Adams, who teaches Information Assurance and stood
at the head of the classroom during the April exercise.

Cadet Joshua Ewing said he and his fellow Blue Team members ``learn all the techniques that a hacker
would do, and we try to beat a hacker.''

These strategies are not just theoretical. Most of these cadets will soon be sent to Afghanistan to
carry out such work, Cadet Ewing said.

When the military deploys in a combat zone or during a domestic emergency, establishing a secure
Internet connection is an early priority. To keep things humming, the military's experts must fend
off the ordinary chaos of the Internet as well as attacks devised to disable the communications
system, like flooding e-mail servers with so many junk messages that they collapse.

Underscoring how seriously the cadets were taking the April games, the sign above the darkened
entranceway in Thayer Hall read ``Information Warfare Live Fire Range'' and the area was draped with
camouflage netting.

One group had to retrieve crucial information from a partly erased hard drive. One common method of
hiding text, said Cadet Sean Storey, is to embed it in digital photographs; he had managed to find
secret documents hidden this way. He was seeking a password needed to read encrypted e-mail he had
located on the hard drive.

Other cadets worked in tandem, as if plugging a leaky dam, to keep the entire system working as the
N.S.A. hackers attacked the engine that runs a crucial database as well as the e-mail server.

They shouted out various Internet addresses to inspect -- and usually block -- after getting
clearance from referees. And there was that awkward moment when the cadet in charge, Salvatore
Messina, had to act without clearance because the attack was so severe he couldn't even send an
e-mail message.

The cadets in this room do get their share of ribbing. But one cadet, Derek Taylor, said today's
soldiers recognize that technological expertise can be as vital as brute force in saving lives. West
Point takes the competition seriously. The cadets who helped install and secure the operating system
spent a week setting it up. The dean gives a pep talk; professors bring food.

Brian McCord, part of the team that installed the operating system, said he was chosen because his
senior project was deeply reliant on Linux. The West Point team used this open-source operating
system, freely available on the Internet, instead of relying on proprietary products from big-name
companies like Microsoft or Sun Microsystems.

``It seems weird for the Army with its large contracts to be using Linux, but it's very cheap and
very customizable,'' Cadet McCord said. It is also much easier to secure because ``you can tweak it
for everything you need'' and there are not as many known ways to attack it, he said.

West Point emerged victorious in the games last month. That means the academy, which has won five of
the last nine competitions, can keep the Director's Cup trophy, which is displayed near a German
Enigma encoding machine from World War II. Cracking the Enigma code helped the Allies win the war,
and the machine is a stark reminder of the pivotal role of technology in warfare.

\section{China Holds an Ex-Leader of '89 Rallies, Family Says}

BEIJING -- A onetime student leader of the 1989 Tiananmen Square pro-democracy movement who was
seized and secretly imprisoned by Chinese authorities as he sought to re-enter China last October
has been charged with fraud, his family said Wednesday.

News of the charge against the former student leader, Zhou Yongjun, who headed the Beijing Students'
Autonomous Union early in the movement, surfaced only weeks before the June 4 anniversary of the
bloody crackdown that killed hundreds of demonstrators and ended the Tiananmen Square movement.

China's government has gone to great lengths to ensure that the 20th anniversary of the crackdown, a
momentous date for many human-rights advocates, passes unnoticed here.

The authorities filed the fraud charge in Suining, the town in Sichuan Province in southwest China
where Mr.~Zhou's parents live and where he is jailed, his older sister, Zhou Shufen, said in a phone
interview from Chengdu, Sichuan's capital. Ms.~Zhou said her family in Suining told her that the
local Public Security Bureau had delivered a notice of the charge to their home.

Neither she nor other family members have been allowed to see Mr.~Zhou. But fellow prisoners and
officials have told them that his case was ``political,'' she said, and that other inmates have been
told to avoid him.

``Our family demands that they release him immediately and call to account those who have illegally
detained him,'' Ms.~Zhou said.

Mr.~Zhou's partner, Zhang Yuewei, said in a telephone interview from their California home that he
had gone back to China to see his parents because they were in poor health.

``He also wanted to take his parents to the States to emigrate, so they could have a better life,''
she said.

John Kusumi, the founder of the China Support Network, an organization in the United States formed
after the 1989 protests, said that Mr.~Zhou apparently was detained as he tried to cross from Hong
Kong into mainland China in early October.

Mr.~Zhou, who is also known by his English name Majer Zhou, fled China in the early 1990s and has
been a permanent resident of the United States since 1993.

A spokesman for the American Embassy in Beijing said that officials there had heard reports of the
arrest and were looking into it but had no other comment.

Mr.~Zhou, now 41, was an undergraduate at the China University of Political Science and Law in
Beijing in 1989, when students began to flock to Tiananmen Square to honor Hu Yaobang, the
reform-minded Chinese leader whose death set off the democracy protests.

Mr.~Zhou was elected the first president of the new Students' Autonomous Union but lost the position
after a disagreement over whether to proceed with a crucial demonstration in late April of that
year. Mr.~Zhou remained part of the democracy movement, however, and was in Tiananmen Square during
the June 4 military assault.

He served two years in prison for his role in the protest, then fled to Hong Kong and, later, the
United States. He tried to re-enter China in 1998 and was imprisoned for three more years.

Also on Wednesday, the Dui Hua Foundation, which advocates for human rights in China, lowered its
estimate of the number of Tiananmen protesters who remained imprisoned.

The foundation, based in San Francisco, said it believed that about 30 people remained in jail, down
from the 50 to 60 prisoners it had previously reported. The group lowered its estimate after Chinese
officials notified it of the early release of several former protesters, and other researchers
documented the status of 104 demonstrators who had been jailed in Beijing No.~2 Prison after the June 4 crackdown. All but six of those prisoners have been released, the foundation stated.

\section{Thriving Norway Provides an Economics Lesson}

When capitalism seemed on the verge of collapse last fall, Kristin Halvorsen, Norway's Socialist
finance minister and a longtime free market skeptic, did more than crow.

As investors the world over sold in a panic, she bucked the tide, authorizing Norway's \$300 billion
sovereign wealth fund to ramp up its stock buying program by \$60 billion -- or about 23 percent of
Norway 's economic output.

``The timing was not that bad,'' Ms.~Halvorsen said, smiling with satisfaction over the broad
worldwide market rally that began in early March.

The global financial crisis has brought low the economies of just about every country on earth. But
not Norway.

With a quirky contrariness as deeply etched in the national character as the fjords carved into its
rugged landscape, Norway has thrived by going its own way. When others splurged, it saved. When
others sought to limit the role of government, Norway strengthened its cradle-to-grave welfare
state.

And in the midst of the worst global downturn since the Depression, Norway's economy grew last year
by just under 3 percent. The government enjoys a budget surplus of 11 percent and its ledger is
entirely free of debt.

By comparison, the United States is expected to chalk up a fiscal deficit this year equal to 12.9
percent of its gross domestic product and push its total debt to \$11 trillion, or 65 percent of the
size of its economy.

Norway is a relatively small country with a largely homogeneous population of 4.6 million and the
advantages of being a major oil exporter. It counted \$68 billion in oil revenue last year as prices
soared to record levels. Even though prices have sharply declined, the government is not
particularly worried. That is because Norway avoided the usual trap that plagues many energy-rich
countries.

Instead of spending its riches lavishly, it passed legislation ensuring that oil revenue went
straight into its sovereign wealth fund, state money that is used to make investments around the
world. Now its sovereign wealth fund is close to being the largest in the world, despite losing 23
percent last year because of investments that declined.

Norway's relative frugality stands in stark contrast to Britain, which spent most of its North Sea
oil revenue -- and more -- during the boom years. Government spending rose to 47 percent of G.D.P.,
from 42 percent in 2003. By comparison, public spending in Norway fell to 40 percent from 48 percent
of G.D.P.

``The U.S.~and the U.K. have no sense of guilt,'' said Anders Aslund, an expert on Scandinavia at
the Peterson Institute for International Economics in Washington. ``But in Norway, there is instead
a sense of virtue. If you are given a lot, you have a responsibility.''

Eirik Wekre, an economist who writes thrillers in his spare time, describes Norwegians' feelings
about debt this way: ``We cannot spend this money now; it would be stealing from future
generations.''

Mr.~Wekre, who paid for his house and car with cash, attributes this broad consensus to as the
country's iconoclasm. ``The strongest man is he who stands alone in the world,'' he said, quoting
Norwegian playwright Henrik Ibsen.

Still, even Ibsen might concede that it is easier to stand alone when your nation has benefited from
oil reserves that make it the third-largest exporter in the world. The money flowing from that black
gold since the early 1970s has prompted even the flintiest of Norwegians to relax and enjoy their
good fortune. The country's G.D.P. per person is \$52,000, behind only Luxembourg among industrial
democracies.

As in much of the rest of the world home prices have soared here, tripling this decade. But there
has been no real estate crash in Norway because there were few mortgage lending excesses. After a 15
percent correction, prices are again on the rise.

Unlike Dublin or Riyadh, Saudi Arabia, where work has stopped on half-built skyscrapers and stilled
cranes dot the skylines, Oslo retains a feeling of modesty reminiscent of a fishing village rather
than a Western capital, with the recently opened \$800 million Opera House one of the few signs of
opulence.

Norwegian banks, said Arne J.~Isachsen, an economist at the Norwegian School of Management, remain
largely healthy and prudent in their lending. Banks represent just 2 percent of the economy and
tight public oversight over their lending practices have kept Norwegian banks from taking on the
risk that brought down their Icelandic counterparts. But they certainly have not closed their doors
to borrowers. Mr.~Isachsen, like many in Norway, has a second home and an open credit line from his
bank, which he recently used to buy a new boat.

Some here worry that while a cabin in the woods and a boat may not approach the excesses seen in New
York or London, oil wealth and the state largesse have corrupted Norway's once-sturdy work ethic.

``This is an oil-for-leisure program,'' said Knut Anton Mork, an economist at Handelsbanken in Oslo.
A recent study, he pointed out, found that Norwegians work the fewest hours of the citizens of any
industrial democracy.

``We have become complacent,'' Mr.~Mork added. ``More and more vacation houses are being built. We
have more holidays than most countries and extremely generous benefits and sick leave policies. Some
day the dream will end.''

But that day is far off. For now, the air is clear, work is plentiful and the government's helping
hand is omnipresent -- even for those on the margins.

Just around the corner from Norway's central bank, for instance, Paul Bruum takes a needle full of
amphetamines and jabs it into his muscular arm. His scabs and sores betray many years as a heroin
addict. He says that the \$1,500 he gets from the government each month is enough to keep him
well-fed and supplied with drugs.

Mr.~Bruum, 32, says he has never had a job, and he admits he is no position to find one. ``I don't
blame anyone,'' he said. ``The Norwegian government has provided for me the best they can.''

To Ms.~Halvorsen, the finance minister, even the underside of the Norwegian dream looks pretty good
compared to the economic nightmares elsewhere.

``As a socialist, I have always said that the market can't regulate itself,'' she said. ``But even I
was surprised how strong the failure was.''


\section{Europe Fines Intel \$1.45 Billion in Antitrust Case}


European regulators have spent close to nine years investigating whether Intel, the world's largest
chip maker, illegally hampered competitors in its dealings with computer makers and retailers. And
on Wednesday, they finally reached a conclusion, punishing Intel with a record fine of 1.06 billion
euros (\$1.45 billion) for abusing its dominant position in the computer chip market.

The European Commission accused Intel of offering computer makers better deals if they would agree
to sell products that relied solely on Intel's processors or delay products that used chips made by
Intel's rival, Advanced Micro Devices. In addition, regulators said Intel paid a retailer to keep
only Intel chips in stock.

While the commission does not want to eliminate volume discounts, which are common in the PC
industry, it ordered Intel to stop placing onerous conditions on the recipients.

But while A.M.D. cheered the ruling as encouraging challengers to Intel's dominance, the history of
antitrust cases against Intel and other big technology companies suggests that little will change.

As the world's largest chip maker and the co-creator of the PC industry, Intel brings vast
technological and business resources to a game whose rules change at Internet speed. A few minutes
in the penalty box, while perhaps costly, will hardly have any impact on Intel's position.

``There is nothing in this ruling that reverses Moore's Law,'' said Paul S.~Otellini, Intel's chief
executive, during a conference call to discuss the ruling in Europe. He was referring to the
observation made famous by Intel's co-founder, Gordon Moore, that the computing muscle of chips
tends to double about every two years, even as prices fall.

Intel and A.M.D. have squabbled in courtrooms for decades, arguing over issues like Intel's control
over the underlying architecture of PC chips and the company's sales tactics. A.M.D. has achieved a
number of victories in these cases and is the main alternative to Intel as the supplier of the
powerful microprocessors that function as the brains of PCs and computer servers.

However, A.M.D.'s overall position in the chip market has remained relatively constant. Intel tends
to claim about 80 percent of total PC and server chip sales, while A.M.D. takes close to 20 percent.

``That really has never shifted much, whether there are antitrust rulings or not,'' said Dan
Hutcheson, the chief executive of VLSI Research, a chip research firm.

A.M.D. argues that the commission's decision will have a big impact. ``With this ruling, the
industry will benefit from an end to Intel's monopoly-inflated pricing and European consumers will
enjoy greater choice, value and innovation,'' said Thomas M.~McCoy, A.M.D.'s executive vice
president for legal affairs.

Through complaints in Europe, Asia and the United States, A.M.D. has spent vast amounts of time and
money raising concerns about Intel's interactions with the companies that produce and sell PCs and
servers, such as Dell and Acer. A.M.D. has accused Intel of essentially paying companies to not use
A.M.D.'s chips in their products and to delay the release of products based on A.M.D.'s chips.

Last year, Korean regulators fined Intel for harming A.M.D.'s business. And now the European
Commission has followed suit. The competition commissioner, Neelie Kroes, said Intel had skewed
competition and denied consumers a choice for chips, with practices that ``undermined innovation.''

Intel plans to appeal the European ruling and has denied including any exclusionary measures in its
contracts with customers or retailers.

Many analysts charge that A.M.D.'s intense focus on trying to alter Intel's interactions with
customers will have little impact on its overall business.

``When you get down to it, whatever Intel is doing that A.M.D. is complaining about probably isn't
making that much of a difference,'' said Linley Gwennap, a chip analyst at the Linley Group. ``The
Intel discounts or whatever are just noise.''

A.M.D.'s most impressive and profitable moments have come when it has read the technology tea leaves
and developed superior products. At times, A.M.D. has beaten Intel on overall PC chip speeds. In
2003, it introduced a server chip called Opteron that shook up the market and for the first time
made the company a serious player in the data centers of large businesses.

In fact, A.M.D.'s successes during the last few years, achieved while European regulators conducted
their investigation, have altered the chip landscape. A.M.D.'s chips now appear in both the PCs and
the servers of most major manufacturers.

``Europe's case is really predicated on the idea that there will be future harm to A.M.D.,'' said
Geoffrey Manne, who has practiced antitrust law and taught at Lewis \& Clark Law School in Portland,
Ore. ``It is really hard to find evidence of that now.''

One of A.M.D.'s major problems has been its inability to produce top-notch products regularly. After
a series of ground-breaking chips, it often ends up shipping products late or has to rework them to
eliminate troublesome bugs. Intel, meanwhile, uses its vast resources to catch up and then surpass
A.M.D.

``A.M.D. is basically operating on a shoestring,'' Mr.~Gwennap said. ``If you have a slip-up, that
can kill your revenue for the next couple of years. But Intel has projects and back-up projects and
back-ups to the back-ups.''

A.M.D. insists that Intel has just as many problems getting new chips out, but can use its heft to
outflank competitors.

``When you are one of the biggest and richest monopolies of the century, it is pretty easy to hide
your sins,'' said Mr.~McCoy of A.M.D. ``They shoot us in the knee and then tell everyone, 'Hey,
look, they don't run fast.' ''

The industry is paying more attention these days to Intel's dealings in markets other than PCs.
Intel is gunning for Nvidia in the graphics-chip market and for Freescale Semiconductor, Qualcomm
and Texas Instruments in cellphone chips.

Ultimately, the European Commission's ruling appears to address yesterday's problems.

``There is a recognition in the U.S.~that this is a problem, but no court has tried to confront it
or change its approach based on this fear that things change too quickly for a remedy to have any
impact,'' said Keith N.~Hylton, a Boston University law professor.

\section{In Year After Quake, China Sealed an Opened Door}

BEIJING -- The earthquake that killed 87,000 people in Sichuan Province in China a year ago this
week was a devastating tragedy that, through all the rubble, offered a few rays of hope.

In a rare burst of volunteerism, more than a million people flocked to the quake zone. The nation's
leaders demonstrated their ability to alleviate human suffering by mobilizing 146,000 troops in a
matter of days. China's biggest corporations donated nearly \$1 billion. Citizens and celebrities
like Jackie Chan, Jet Li and Yao Ming gave hundreds of millions more.

``The great task of earthquake rescue and recovery reminds us again that unity is strength, and that
victory can only be gained through struggle,'' President Hu Jintao said Tuesday during a somber
commemoration, shown live on national television, observing the anniversary of China's worst natural
disaster in more than three decades.

But if China's response to the earthquake has been widely commended, critics say it did not
ultimately change the government's opaque and authoritarian ways. The relaxed press freedom that
accompanied the first weeks of the recovery effort proved to be fleeting. And after vowing to
thoroughly investigate why so many school buildings crumbled, killing about 5,300 students by
official estimates, the government insisted last week that the wrath of nature, not human
malfeasance, was to blame. ``Up until now, we haven't found that anybody caused or did anything to
make buildings more vulnerable,'' said Tang Kai, who directs the planning office for the Ministry of
Housing and Urban-Rural Development.

Over the past year, as they tried to expose instances of substandard school construction, bloggers,
activists and parents whose children perished under mounds of rubble have been harassed or jailed.
Most parents acknowledge that they accepted payments from the government that require them to stay
silent on the matter.

In recent weeks, as the anniversary approached, scores of foreign journalists trying to interview
aggrieved parents were intimidated and roughed up.

Ai Weiwei, an artist and architect who has organized a campaign to count the number of dead
children, said that 20 of his volunteers had been attacked or detained and that his blog postings
were frequently deleted. ``What the government has done is irresponsible,'' said Mr.~Ai, one of the
few critics who has remained unscathed. ``The Chinese people deserve better.''

In the early days of the disaster, there were reasons for encouragement. Hours after the quake
struck, Prime Minister Wen Jiabao flew into the quake zone trailed by television cameras. In a
country with a history of concealing bad news, Mr.~Wen's very public role consoling survivors and
shouting encouraging words to those trapped in the ruins of a flattened school seemed to signal that
this disaster would be handled differently.

``At the beginning, Chinese journalists were very excited to be covering a major natural disaster in
what were the most free conditions ever,'' said Qian Gang, the author of ``The Great China
Earthquake,'' an account of the 1976 quake that devastated the city of Tangshan and killed at least
240,000 people, which was initially covered up. ``But gradually the power of the government seeped
back in.''

In an article published last week by China Media Project, a Web site run by the University of Hong
Kong, Mr.~Qian detailed how many journalists, taking advantage of the chaos, ignored reporting
limitations imposed by the central propaganda department. Even dispatches from Xinhua, the state-run
news agency, noted that collapsed schools were often surrounded by buildings that appeared
unscathed.

Four days after the earthquake, an online chat sponsored by People's Daily, the Communist Party's
mouthpiece, featured officials who stated that ``there were certainly quality issues behind the
collapse of school buildings, and we will conduct a strict and uncompromising investigation.''

By late June, however, the reins had been fully pulled back. Reporters who did not work for Xinhua
or other strictly controlled, state-run media outlets were forced to leave the quake zone.
Publications that had questioned the quality of the school buildings were criticized, or in some
cases compelled to punish the offending writers and editors.

One reporter said that he and his colleagues had gathered compelling proof that cutting corners and
inadequate oversight by local officials contributed to the high student death toll. But, he said, it
has become impossible to publish those findings. ``At this point, stability outweighs everything
else,'' said the reporter, who spoke on the condition of anonymity and asked that the name of his
employer not be disclosed.

Officials have been less successful in silencing the scores of parents who continue to press for an
official accounting of why so many schools collapsed that day.

Liu Xiaoying of Mianzhu, speaking by telephone, explained how she and other parents were not allowed
to join the public mourning on Tuesday for fear that they might disrupt the tightly orchestrated
commemorative events that were replete with foreign dignitaries. A phalanx of armed officers made
sure reporters did not venture into town.

``There has been no government investigation report at all,'' said Ms.~Liu, whose daughter died at
Fuxin No.~2 Primary School. ``If there had been, perhaps our grief would not be this great.''

\pagebreak
\printindex

\end{document}
